\documentclass[aspectratio=169]{beamer}

% ── Theme & Colors ────────────────────────────────────────────────────────────
\usetheme{Madrid}
\usecolortheme{default}

\definecolor{primary}{RGB}{30, 60, 114}
\definecolor{accent}{RGB}{0, 150, 199}
\definecolor{lightgray}{RGB}{245, 245, 248}
\definecolor{successgreen}{RGB}{39, 174, 96}
\definecolor{warnred}{RGB}{192, 57, 43}
\definecolor{warnyellow}{RGB}{241, 196, 15}

\setbeamercolor{palette primary}{bg=primary, fg=white}
\setbeamercolor{palette secondary}{bg=accent, fg=white}
\setbeamercolor{palette tertiary}{bg=primary!80, fg=white}
\setbeamercolor{palette quaternary}{bg=primary, fg=white}
\setbeamercolor{structure}{fg=primary}
\setbeamercolor{section in toc}{fg=primary}
\setbeamercolor{block title}{bg=primary, fg=white}
\setbeamercolor{block body}{bg=lightgray, fg=black}
\setbeamercolor{alertblock title}{bg=warnred, fg=white}
\setbeamercolor{alertblock body}{bg=warnred!10, fg=black}
\setbeamercolor{exampleblock title}{bg=successgreen, fg=white}
\setbeamercolor{exampleblock body}{bg=successgreen!10, fg=black}

% ── Packages ──────────────────────────────────────────────────────────────────
\usepackage{tikz}
\usetikzlibrary{arrows.meta, positioning, shapes.geometric, fit, backgrounds}
\usepackage{fontawesome5}
\usepackage{booktabs}
\usepackage{array}
\usepackage{xcolor}
\usepackage{graphicx}
\usepackage{hyperref}

% ── Metadata ──────────────────────────────────────────────────────────────────
\title[Agentic RAG — Weekly Report]{\textbf{Agentic RAG Implementation}\\Weekly Progress Report}
\subtitle{Challenges, Architecture \& Roadmap}
\author[Brahim \& Youcef]{\textbf{Brahim} \and \textbf{Youcef}}
\date{\today}
\institute{Research \& Development Team}

% ── Helpers ───────────────────────────────────────────────────────────────────
\newcommand{\pro}[1]{\textcolor{successgreen}{\faCheckCircle\; #1}}
\newcommand{\con}[1]{\textcolor{warnred}{\faTimesCircle\; #1}}
\newcommand{\todo}[1]{\textcolor{accent}{\faArrowCircleRight\; #1}}

% ─────────────────────────────────────────────────────────────────────────────
\begin{document}

% ── Title Slide ──────────────────────────────────────────────────────────────
{
\setbeamertemplate{headline}{}
\begin{frame}[plain]
  \vspace{1.2cm}
  \begin{center}
    {\usebeamercolor[fg]{structure}\Huge\textbf{Agentic RAG Implementation}}\\[0.3cm]
    {\large\textit{Weekly Progress Report}}\\[0.5cm]
    \rule{8cm}{0.4pt}\\[0.4cm]
    {\large \faUser\ \textbf{Brahim} \quad\&\quad \faUser\ \textbf{Youcef}}\\[0.3cm]
    {\small \today}
  \end{center}
  \vfill
  \begin{block}{}
    \centering\small
    Based on the architecture proposed in \emph{``Developing an AI Assistant for Knowledge Management\\
    and Workforce Training in State DOTs''} — Amaram et al., arXiv:2603.03302v1 (2026)
  \end{block}
\end{frame}
}

% ── Outline ──────────────────────────────────────────────────────────────────
\begin{frame}{Agenda}
  \tableofcontents
\end{frame}

% ─────────────────────────────────────────────────────────────────────────────
\section{Background — The Paper \& Its Architecture}
% ─────────────────────────────────────────────────────────────────────────────

\begin{frame}{Reference Paper}
  \begin{columns}[T]
    \begin{column}{0.55\textwidth}
      \begin{block}{Paper Overview}
        \small
        \begin{itemize}
          \item \textbf{Title:} Developing an AI Assistant for Knowledge Management and Workforce Training in State DOTs
          \item \textbf{Authors:} Amaram et al. (University of Houston)
          \item \textbf{Source:} arXiv:2603.03302v1, Feb 2026
        \end{itemize}
      \end{block}
      \vspace{0.3cm}
      \begin{exampleblock}{Key Results}
        \small
        \begin{itemize}
          \item \textbf{Corpus:} 521 technical DOT documents
          \item \textbf{Evaluation:} 100 domain-specific queries
          \item \textbf{Precision@3:} $\approx 1.0$
          \item \textbf{Recall@3:} $94.4\%$
          \item Single-pass RAG Recall@5: only $58\%$
        \end{itemize}
      \end{exampleblock}
    \end{column}
    \begin{column}{0.42\textwidth}
      \begin{alertblock}{Why We Chose This Paper}
        \small
        \begin{itemize}
          \item Clear, modular multi-agent design
          \item Reproducible agent prompts provided
          \item Strong improvement over standard RAG
          \item Framework-agnostic architecture
        \end{itemize}
      \end{alertblock}
      \vspace{0.3cm}
      \begin{block}{Our Goal}
        \small
        Reimplement \textbf{the architecture} (not the dataset or domain) to understand and extend Agentic RAG.
      \end{block}
    \end{column}
  \end{columns}
\end{frame}

% ── Architecture ─────────────────────────────────────────────────────────────
\begin{frame}{The Multi-Agent RAG Architecture}
  \begin{center}
  \begin{tikzpicture}[
    node distance=0.7cm and 0.55cm,
    agent/.style={
      rectangle, rounded corners=4pt,
      minimum width=2.1cm, minimum height=0.75cm,
      draw=primary, fill=primary!10,
      font=\small\bfseries, align=center, text=primary
    },
    arrow/.style={-{Stealth[length=5pt]}, thick, color=accent},
    label/.style={font=\scriptsize, color=gray}
  ]
    \node[agent] (user)  {\faUser\\User Proxy\\Agent};
    \node[agent, right=of user]  (ret)  {\faSearch\\Retriever\\Agent};
    \node[agent, right=of ret]   (gen)  {\faEdit\\Generator\\Agent};
    \node[agent, right=of gen]   (eval) {\faCheckSquare\\Evaluator\\Agent};
    \node[agent, below=0.9cm of eval] (ref)  {\faSyncAlt\\Query Refiner\\Agent};

    \draw[arrow] (user) -- node[above, label]{query} (ret);
    \draw[arrow] (ret)  -- node[above, label]{context} (gen);
    \draw[arrow] (gen)  -- node[above, label]{draft} (eval);
    \draw[arrow] (eval) -- node[right, label]{feedback} (ref);
    \draw[arrow] (ref.west) -- ++(-0.4,0) |- node[left, label, pos=0.3]{refined query} (ret.south);

    \node[below=0.15cm of user,  font=\tiny, color=gray] {\circled{1}};
    \node[below=0.15cm of ret,   font=\tiny, color=gray] {\circled{2}};
    \node[below=0.15cm of gen,   font=\tiny, color=gray] {\circled{3}};
    \node[below=0.15cm of eval,  font=\tiny, color=gray] {\circled{4}};
    \node[below=0.15cm of ref,   font=\tiny, color=gray] {\circled{5}};

    % Satisfactory exit
    \draw[arrow, color=successgreen] (eval.north) -- ++(0,0.55) node[right, font=\scriptsize, color=successgreen]{\faCheckCircle\ Final Answer};
  \end{tikzpicture}
  \end{center}

  \vspace{0.2cm}
  \begin{columns}[T]
    \small
    \begin{column}{0.19\textwidth}
      \textbf{\textcolor{primary}{\circled{1} User Proxy}}\\
      Receives \& relays query without modification
    \end{column}
    \begin{column}{0.19\textwidth}
      \textbf{\textcolor{primary}{\circled{2} Retriever}}\\
      Similarity search over vector DB
    \end{column}
    \begin{column}{0.19\textwidth}
      \textbf{\textcolor{primary}{\circled{3} Generator}}\\
      Formats \& cites a grounded answer
    \end{column}
    \begin{column}{0.19\textwidth}
      \textbf{\textcolor{primary}{\circled{4} Evaluator}}\\
      Checks clarity, relevance, completeness
    \end{column}
    \begin{column}{0.19\textwidth}
      \textbf{\textcolor{primary}{\circled{5} Refiner}}\\
      Rewrites query on bad feedback (max $k$ times)
    \end{column}
  \end{columns}
\end{frame}

% ─────────────────────────────────────────────────────────────────────────────
\section{Week 1 — Langflow Experiment}
% ─────────────────────────────────────────────────────────────────────────────

\begin{frame}{What Is Langflow?}
  \begin{columns}[T]
    \begin{column}{0.5\textwidth}
      \begin{block}{Overview}
        \small
        Langflow is a \textbf{low-code, visual} framework for building LLM-powered pipelines. It offers a drag-and-drop interface to connect components (loaders, retrievers, LLMs, agents) without writing boilerplate code.
      \end{block}
      \vspace{0.3cm}
      \begin{block}{Why We Started With It}
        \small
        \begin{itemize}
          \item Fast prototyping with visual canvas
          \item Built-in RAG components
          \item No infrastructure setup needed
          \item Seemed ideal for a quick proof-of-concept
        \end{itemize}
      \end{block}
    \end{column}
    \begin{column}{0.47\textwidth}
      \begin{alertblock}{\faExclamationTriangle\ Reality Check}
        \small
        Despite the promise of simplicity, we quickly ran into fundamental limitations when trying to replicate a \textbf{multi-agent} workflow rather than a simple single-chain RAG pipeline.
      \end{alertblock}
    \end{column}
  \end{columns}
\end{frame}

\begin{frame}{Phase 1 — Initial Langflow Attempt}
  \begin{block}{What We Tried}
    \small Reproduce the paper's full 5-agent pipeline (User Proxy $\to$ Retriever $\to$ Generator $\to$ Evaluator $\to$ Query Refiner) inside Langflow's visual canvas.
  \end{block}
  \vspace{0.3cm}
  \begin{columns}[T]
    \begin{column}{0.48\textwidth}
      \begin{alertblock}{\faClock\ Time Lost}
        \small
        \begin{itemize}
          \item Spent significant time navigating Langflow's documentation and community forums
          \item Repeated attempts to wire conditional feedback loops between agents
          \item Iterating on workarounds that ultimately failed
        \end{itemize}
      \end{alertblock}
    \end{column}
    \begin{column}{0.48\textwidth}
      \begin{alertblock}{\faBan\ Core Blockers}
        \small
        \begin{itemize}
          \item \textbf{No native multi-agent support:} Langflow is built around sequential chains, not autonomous agent graphs
          \item \textbf{No conditional loops:} Cannot model ``if unsatisfactory $\to$ refine \& retry''
          \item \textbf{Limited agent customization:} Agent prompts and behaviors are constrained by UI components
        \end{itemize}
      \end{alertblock}
    \end{column}
  \end{columns}
\end{frame}

\begin{frame}{What We Managed to Build in Langflow}
  \begin{columns}[T]
    \begin{column}{0.48\textwidth}
      \begin{exampleblock}{\faCheck\ Partial Successes}
        \small
        \begin{itemize}
          \item Basic document ingestion and chunking pipeline
          \item Vector embedding and storage with a supported DB ( AstraDB )
          \item Simple single-pass retrieval + generation chain
          \item Verified that the RAG concept works end-to-end at a basic level
        \end{itemize}
      \end{exampleblock}
    \end{column}
    \begin{column}{0.48\textwidth}
      \begin{alertblock}{\faTimes\ What We Could Not Implement}
        \small
        \begin{itemize}
          \item Evaluator Agent with structured pass/fail logic
          \item Query Refiner Agent feeding back into retrieval
          \item Iterative refinement loop (max $k$ retries)
          \item Coordination between more than 2 agents
          \item Figure preprocessing via a VLM
        \end{itemize}
      \end{alertblock}
    \end{column}
  \end{columns}
  \vspace{0.4cm}
  \begin{block}{Key Takeaway}
    \centering\small
    Langflow is suited for \textbf{linear RAG pipelines}. Implementing a \textbf{graph-based, stateful, multi-agent} workflow requires a different paradigm entirely.
  \end{block}
\end{frame}

% ─────────────────────────────────────────────────────────────────────────────
\section{Gap Analysis — Paper vs. Our Implementation}
% ─────────────────────────────────────────────────────────────────────────────

\begin{frame}{Gap Analysis}
  \begin{center}
  \renewcommand{\arraystretch}{1.35}
  \begin{tabular}{>{\bfseries\small}l >{\small}c >{\small}c}
    \toprule
    \textbf{Component} & \textbf{Paper} & \textbf{Our Langflow Build} \\
    \midrule
    Document ingestion \& chunking     & \textcolor{successgreen}{\faCheckCircle} & \textcolor{successgreen}{\faCheckCircle} \\
    Vector embedding \& storage        & \textcolor{successgreen}{\faCheckCircle} & \textcolor{successgreen}{\faCheckCircle} \\
    Single-pass retrieval + generation & \textcolor{successgreen}{\faCheckCircle} & \textcolor{successgreen}{\faCheckCircle} \\
    User Proxy Agent                   & \textcolor{successgreen}{\faCheckCircle} & \textcolor{warnyellow}{\faMinus} (simulated) \\
    Retriever Agent                    & \textcolor{successgreen}{\faCheckCircle} & \textcolor{warnyellow}{\faMinus} (partial) \\
    Generator Agent                    & \textcolor{successgreen}{\faCheckCircle} & \textcolor{warnyellow}{\faMinus} (partial) \\
    Evaluator Agent                    & \textcolor{successgreen}{\faCheckCircle} & \textcolor{warnred}{\faTimesCircle} \\
    Query Refiner Agent                & \textcolor{successgreen}{\faCheckCircle} & \textcolor{warnred}{\faTimesCircle} \\
    Iterative refinement loop          & \textcolor{successgreen}{\faCheckCircle} & \textcolor{warnred}{\faTimesCircle} \\
    Figure preprocessing (VLM)         & \textcolor{successgreen}{\faCheckCircle} & \textcolor{warnred}{\faTimesCircle} \\
    \bottomrule
  \end{tabular}
  \end{center}
  \vspace{0.2cm}
  \small\centering
  \textcolor{successgreen}{\faCheckCircle}\ Implemented \quad
  \textcolor{warnyellow}{\faMinus}\ Partial \quad
  \textcolor{warnred}{\faTimesCircle}\ Not implemented
\end{frame}

% ─────────────────────────────────────────────────────────────────────────────
\section{Next Steps — Changing the Development Stack}
% ─────────────────────────────────────────────────────────────────────────────

\begin{frame}{Why We Need a New Stack}
  \begin{columns}[T]
    \begin{column}{0.48\textwidth}
      \begin{block}{Requirements for Agentic RAG}
        \small
        \begin{itemize}
          \item \textbf{Graph-based execution:} agents as nodes with conditional edges
          \item \textbf{Stateful memory:} pass context, feedback, and history between agents
          \item \textbf{Custom agent logic:} full control over prompts and tool calls
          \item \textbf{Conditional loops:} ``retry if unsatisfactory'' with a max iteration cap
          \item \textbf{Modularity:} easy to swap LLMs, vector DBs, or embedding models
          \item \textbf{Observability:} trace every agent step for debugging
        \end{itemize}
      \end{block}
    \end{column}
    \begin{column}{0.48\textwidth}
      \begin{alertblock}{Langflow Verdict}
        \small
        Langflow \textbf{does not satisfy} these requirements. Its visual pipeline model is inherently linear and lacks the primitives needed for graph-based agent orchestration.
      \end{alertblock}
      \vspace{0.3cm}
      \begin{block}{Decision}
        \small
        We will migrate to a \textbf{code-first framework} that provides native support for multi-agent coordination, stateful graphs, and flexible tool use.
      \end{block}
    \end{column}
  \end{columns}
\end{frame}

\begin{frame}{Candidate Stacks — Overview}
  \begin{center}
  \renewcommand{\arraystretch}{1.3}
  \footnotesize
  \begin{tabular}{>{\bfseries}l l l l l}
    \toprule
    \textbf{Framework} & \textbf{Paradigm} & \textbf{Multi-Agent} & \textbf{Conditional Loops} & \textbf{Complexity} \\
    \midrule
    LangGraph          & Graph / DAG       & \textcolor{successgreen}{\faCheckCircle} Native  & \textcolor{successgreen}{\faCheckCircle} Yes & Medium \\
    AutoGen            & Conversation      & \textcolor{successgreen}{\faCheckCircle} Native  & \textcolor{successgreen}{\faCheckCircle} Yes & Medium \\
    CrewAI             & Role-based        & \textcolor{successgreen}{\faCheckCircle} Native  & \textcolor{successgreen}{\faCheckCircle} Yes & Low–Medium \\
    LlamaIndex         & RAG-first         & \textcolor{successgreen}{\faCheckCircle} Native  & \textcolor{successgreen}{\faCheckCircle} Yes & Medium \\
    Haystack           & Pipeline          & \textcolor{warnyellow}{\faMinus} Partial & \textcolor{warnyellow}{\faMinus} Limited & Low \\
    Semantic Kernel    & Plugin-based      & \textcolor{successgreen}{\faCheckCircle} Native  & \textcolor{successgreen}{\faCheckCircle} Yes & Medium–High \\
    Custom Python      & Free-form         & \textcolor{successgreen}{\faCheckCircle} Full    & \textcolor{successgreen}{\faCheckCircle} Full & High \\
    \bottomrule
  \end{tabular}
  \end{center}
\end{frame}

\begin{frame}{Stack Deep-Dive (1/2)}
  \begin{columns}[T]
    \begin{column}{0.48\textwidth}
      \begin{exampleblock}{\faProjectDiagram\ LangGraph (LangChain)}
        \small
        \begin{itemize}
          \item Models agents as \textbf{nodes} and decisions as \textbf{conditional edges} in a directed graph
          \item Native support for stateful memory across nodes
          \item Built-in streaming, checkpointing, and human-in-the-loop
          \item Large ecosystem: hundreds of LLM/tool integrations
          \item \textbf{Best fit} for reproducing the paper's feedback loop
        \end{itemize}
      \end{exampleblock}
    \end{column}
    \begin{column}{0.48\textwidth}
      \begin{exampleblock}{\faRobot\ AutoGen (Microsoft Research)}
        \small
        \begin{itemize}
          \item \textbf{Conversation-centric} multi-agent framework (used in the paper itself)
          \item Agents communicate via structured messages
          \item Supports ``GroupChat'' for orchestrating $N$ agents
          \item Built-in code execution and tool-calling agents
          \item Strong academic backing; close match to the paper's original implementation
        \end{itemize}
      \end{exampleblock}
    \end{column}
  \end{columns}
  \vspace{0.3cm}
  \begin{columns}[T]
    \begin{column}{0.48\textwidth}
      
     
    \end{column}
    \begin{column}{0.48\textwidth}
      
    \end{column}
  \end{columns}
\end{frame}

\begin{frame}{Stack Deep-Dive (2/2)}
  \begin{columns}[T]
    \begin{column}{0.48\textwidth}
      \vspace{0.3cm}
      \begin{block}{\faMicrosoft\ Semantic Kernel}
        \small
        \begin{itemize}
          \item Microsoft's enterprise-focused AI orchestration SDK
          \item Plugin architecture for tool integration
          \item Multi-agent support via Agent Framework
          \item Best suited for .NET / Azure-native environments
        \end{itemize}
      \end{block}
            \begin{alertblock}{\faStar\ Our Recommendation}
        \small
        \textbf{LangGraph} or \textbf{AutoGen} — both support the exact graph + feedback-loop structure the paper describes. AutoGen has the advantage of being the original framework used by the authors.
      \end{alertblock}

    \end{column}
    \begin{column}{0.48\textwidth}
      \begin{block}{\faPython\ Custom Python}
        \small
        \begin{itemize}
          \item Full control over every aspect of agent logic
          \item Use any LLM API (OpenAI, Anthropic, Ollama, HuggingFace)
          \item Pair with ChromaDB / Qdrant / Weaviate for vector storage
          \item Use \texttt{sentence-transformers} for embeddings
          \item Higher development effort but zero black-box behavior
        \end{itemize}
      \end{block}
      \vspace{0.3cm}
    \end{column}
  \end{columns}
\end{frame}

% ─────────────────────────────────────────────────────────────────────────────
\section{Summary \& Action Plan}
% ─────────────────────────────────────────────────────────────────────────────

\begin{frame}{Summary}
  \begin{columns}[T]
    \begin{column}{0.48\textwidth}
      \begin{alertblock}{\faHistory\ What Happened This Week}
        \small
        \begin{enumerate}
          \item Studied the multi-agent RAG architecture from the paper
          \item Attempted full 5-agent implementation in \textbf{Langflow}
          \item Lost significant time debugging Langflow's limitations
          \item Achieved a basic single-pass RAG chain only
          \item Concluded Langflow is \textbf{not suitable} for this architecture
        \end{enumerate}
      \end{alertblock}
    \end{column}
    \begin{column}{0.48\textwidth}
      \begin{exampleblock}{\faRoad\ What's Next}
        \small
        \begin{enumerate}
          \item \textbf{Choose} a new framework (LangGraph or AutoGen)
          \item \textbf{Set up} dev environment with chosen stack
          \item \textbf{Implement} all 5 agents with proper state management
          \item \textbf{Add} iterative evaluation--refinement loop
          \item \textbf{Test} on a small document corpus
          \item \textbf{Evaluate} with Precision@k and Recall@k
        \end{enumerate}
      \end{exampleblock}
    \end{column}
  \end{columns}
  \vspace{0.4cm}
  \begin{block}{}
    \centering\small
    \faLightbulb\quad The week was not wasted — we now have a clear picture of \textbf{what does not work} and \textbf{exactly why}, which is a necessary step toward choosing the right tool.
  \end{block}
\end{frame}

\begin{frame}{Action Plan}
  \begin{center}
  \renewcommand{\arraystretch}{1.4}
  \small
  \begin{tabular}{>{\bfseries}c l >{\bfseries}l l}
    \toprule
    \textbf{\#} & \textbf{Task} & & \\
    \midrule
    1 & Finalize framework choice (LangGraph vs.\ AutoGen) & & \\
    2 & Set up project structure \& environment             & & \\
    3 & Implement Retriever Agent + vector DB               & & \\
    4 & Implement Generator Agent with prompt template      & & \\
    5 & Implement Evaluator Agent + scoring logic           & & \\
    6 & Implement Query Refiner + feedback loop             & & \\
    7 & End-to-end integration \& testing                  & & \\
    8 & Measure Precision@3 \& Recall@3                    & & \\
    \bottomrule
  \end{tabular}
  \end{center}
\end{frame}

% ── Thank You ────────────────────────────────────────────────────────────────
{
\setbeamertemplate{headline}{}
\begin{frame}[plain]
  \vfill
  \begin{center}
    {\Huge\textbf{\textcolor{primary}{Thank You}}}\\[0.5cm]
    \rule{6cm}{0.4pt}\\[0.4cm]
    {\large Questions \& Discussion}\\[0.8cm]
    \faUser\ \textbf{Brahim} \qquad \faUser\ \textbf{Youcef}
  \end{center}
  \vfill
  \begin{block}{}
    \centering\footnotesize
    Reference: Amaram et al., \textit{``Developing an AI Assistant for Knowledge Management and Workforce Training in State DOTs''}, arXiv:2603.03302v1, February 2026.
  \end{block}
\end{frame}
}

\end{document}